\chapter{Endovenous Ablation}

\section*{Definition and Overview}
Endovenous ablation is a minimally invasive treatment for varicose veins and other problem veins in the legs. A thin catheter is inserted into the vein, and heat energy, in the form of laser or radiofrequency, is used to close the vein from the inside. Once closed, the vein is no longer used, and blood is redirected to healthy veins. The procedure is performed through a small puncture, usually with local anaesthetic, and most people return to normal activities within a few days. It has largely replaced traditional vein stripping surgery.

\section*{Epidemiology}
Varicose veins are very common, affecting a significant proportion of adults, more often women, and becoming more common with age and pregnancy. Not all varicose veins need treatment, but many people seek treatment for symptoms such as aching, swelling, and itching, and for cosmetic reasons. Endovenous ablation has become the standard treatment for the larger veins that cause symptoms, with excellent results.

\section*{Aetiology and Pathophysiology}
Varicose veins result from the failure of the one-way valves in the leg veins.
\begin{itemize}
    \item \textbf{Valve failure:} the valves that prevent blood flowing backwards in the veins stop working, so blood pools in the superficial veins
    \item \textbf{Pressure:} the pooled blood stretches the veins, making them twisted, bulging, and painful
    \item \textbf{Risk factors:} family history, female sex, pregnancy, age, obesity, and standing occupations
    \item \textbf{The treatment principle:} closing the abnormal vein removes it from the circulation, and blood returns through the deep veins
    \item \textbf{How ablation works:} heat damages the vein wall, which collapses and seals shut; the body then absorbs the vein over time
    \item \textbf{Effects:} symptoms improve, and blood flow to the legs is maintained through healthy veins
\end{itemize}

\section*{Clinical Features}
People seek treatment for varicose veins with particular symptoms.
\begin{itemize}
    \item Aching, heaviness, and tiredness in the legs, worse at the end of the day
    \item Swelling of the ankles and legs
    \item Itching and burning over the veins
    \item Visible, bulging, twisted veins
    \item Skin changes: discolouration, eczema, and hardening around the ankle
    \item Ulcers in advanced venous disease
    \item Cosmetic concerns
\end{itemize}

\section*{Differential Diagnosis}
Other conditions cause leg swelling and pain.
\begin{itemize}
    \item Deep vein thrombosis, a blood clot in the deep veins, which is an emergency
    \item Chronic venous insufficiency without varicose veins
    \item Lymphoedema
    \item Heart, kidney, and liver disease causing swelling
    \item Muscle and joint pain in the legs
    \item Peripheral artery disease
\end{itemize}

\section*{When to Seek Medical Care}
People with symptomatic varicose veins should be assessed by a GP, who can arrange a vascular assessment and referral. Treatment is considered when symptoms are troublesome, complications develop, or the veins are causing skin damage.

\textbf{Call 000 (or local emergency services) for:}
\begin{itemize}
    \item Sudden painful swelling of one leg, which may indicate a blood clot
    \item A bleeding varicose vein
    \item Chest pain or breathlessness, which may indicate a clot travelling to the lungs
    \item Sudden severe leg pain or colour change
\end{itemize}

\section*{Diagnosis and Investigations}
Before treatment, the veins are assessed in detail.
\begin{itemize}
    \item \textbf{Duplex ultrasound:} the key test, which maps the veins, shows the valve function, and identifies the veins to be treated
    \item \textbf{Clinical assessment:} the pattern of symptoms and any skin changes
    \item \textbf{Exclusion of other causes:} assessment for deep vein thrombosis and other conditions
    \item \textbf{Pre-procedure assessment:} general health and any medicines such as blood thinners
\end{itemize}

\section*{Management}
Endovenous ablation is performed as an outpatient procedure.
\begin{itemize}
    \item \textbf{Local anaesthetic:} the leg is numbed with local anaesthetic, often with light sedation
    \item \textbf{The procedure:} a catheter is guided into the vein using ultrasound, and heat is applied to close it; it takes about 45 minutes
    \item \textbf{No stitches:} a small puncture site, closed with a dressing
    \item \textbf{Compression:} stockings or bandages are worn for a period afterwards
    \item \textbf{Recovery:} most people walk immediately and return to work within a few days
    \item \textbf{Combined treatments:} foam sclerotherapy may be used for smaller veins at the same time
    \item \textbf{Follow-up:} ultrasound review to confirm the veins are closed
\end{itemize}

\section*{Complications}
\begin{itemize}
    \item Bruising, tenderness, and numbness around the treated vein, which usually settle
    \item Skin burns or pigmentation changes
    \item Blood clots in the deep veins, which are uncommon but serious
    \item Incomplete closure requiring repeat treatment
    \item Infection, which is rare
    \item Nerve irritation, causing altered sensation
\end{itemize}

\section*{Prognosis and Outcomes}
Endovenous ablation is highly effective, with most treated veins remaining closed and symptoms improving significantly. It is associated with less pain, faster recovery, and fewer complications than traditional surgery. Most people are very satisfied with the results, although new varicose veins can develop over time.

\section*{Prevention and Self-Care}
\begin{itemize}
    \item Wear compression stockings during and after treatment as advised
    \item Walk regularly after the procedure and avoid prolonged standing
    \item Maintain a healthy weight
    \item Exercise regularly, which improves circulation
    \item Elevate the legs when resting
    \item Attend the follow-up ultrasound
    \item Seek prompt care for any leg pain or swelling after the procedure
\end{itemize}

\section*{Sources and Further Reading}
The following reputable sources provide further information for healthcare students and patients seeking reliable, current advice about endovenous ablation.
\begin{itemize}
    \item Australian and New Zealand Society for Vascular Surgery. \textit{Varicose vein treatment information.} https://anzsvs.org.au/
    \item Healthdirect Australia. \textit{Varicose veins and their treatment.} https://www.healthdirect.gov.au/varicose-veins
    \item NHS. \textit{Endovenous laser and radiofrequency ablation.} https://www.nhs.uk/conditions/varicose-veins/
    \item Mayo Clinic. \textit{Varicose veins: treatment options.} https://www.mayoclinic.org/diseases-conditions/varicose-veins/
    \item Better Health Channel. \textit{Varicose veins.} https://www.betterhealth.vic.gov.au/
    \item MedlinePlus. \textit{Endovenous ablation.} https://medlineplus.gov/ency/article/007409.htm
\end{itemize}
